\section{O-ISAC FOUNDATIONS AND COMPARISON FRAMEWORK}
\label{sec:foundations_comparison}
The scope defined in Section~\ref{sec:introduction} brings diverse optical
systems into one survey. This section turns that scope into a comparison
profile by following the physical context, coupling location, measurement
contract, and evidence trace of each result. Together, these elements preserve
the meaning of an observation and show how it can inform the synthesis.
Table~\ref{tab:comparison_record} condenses them into a map for the reader.

% ============================================================================
% FUTURE FIGURE 1 PRODUCTION SPECIFICATION --- DO NOT MATERIALIZE IN THIS PASS
% Intended placement: here, at the opening of the foundations section. It
% bridges the prior-survey positioning in Table I to the comparison vocabulary
% developed below.
%
% Working title:
% Native Evidence Objects in O-ISAC Synthesis
% Stable label: fig:native_evidence_objects
% Blueprint ID: FIG-OISAC-RS1
% Data authority: Table I and the audited 38-source related-synthesis register;
% conceptual illustration with no prevalence denominator.
%
% Reader question:
% Why can results gathered under the O-ISAC label not be placed directly on
% one common numerical scale?
%
% Layout: three equal, parallel, and unconnected horizontal panels. The panels
% are illustrative evidence objects, not stages, families, or quality levels.
%
% Panel A --- model-defined result:
% Assumptions and a system model lead to an analytical bound or optimized
% relation. Attach compact labels for assumptions, objective, and constraints.
%
% Panel B --- medium-native observation:
% A propagation path, physical observable, and measurement setting lead to a
% reported observation. Attach labels for medium, measurement plane, and
% operating conditions.
%
% Panel C --- architecture or deployment object:
% Components and a network or application setting lead to a design or
% deployment relation. Attach labels for shared components, scenario, and
% validation setting.
%
% Common lower band:
% ``The O-ISAC label alone does not supply a common denominator.''
%
% Main visual message:
% The evidence objects differ in meaning, not in rank. Each can be informative
% within its native frame, but its definition and conditions must be retained
% before a cross-study relationship is considered.
%
% Future lead-in (activate only after the figure has been produced and audited):
% Figure~\ref{fig:native_evidence_objects} illustrates the three native
% evidence objects that recur across the survey literature and motivates the
% contextual record developed in this section.
%
% Future caption:
% Native evidence objects represented in O-ISAC syntheses. A model-defined
% result is bounded by analytical assumptions; a medium-native observation is
% tied to its propagation path, observable, and measurement setting; and an
% architecture or deployment object organizes system, network, or application
% choices. The parallel panels are illustrative rather than exhaustive and
% imply neither sequence nor quality. Reader-task families and source
% assignments remain in Table I.
%
% Visual constraints:
% - Do not repeat the six reader-task families or individual sources from
%   Table I; the wider contextual-synthesis register remains an internal audit
%   artifact and is not a reader-facing annex.
% - Do not connect the panels with arrows or arrange them as a sequence,
%   maturity scale, or hierarchy.
% - Do not reproduce the comparison fields or admissibility decisions reserved
%   for Table II and Figure 2.
% - Do not use checkmarks, scores, unequal panel sizes, source counts, or any
%   encoding that could imply prevalence or quality.
% - Use equal panel dimensions, direct labels, and grayscale-safe borders or
%   icons.
% ============================================================================

\subsection{System Boundary, Modality, and Signal Path}

Technical comparison begins by locating communication and sensing within their
physical signal paths. The modality of the observable is recorded separately
from the hardware used to generate, distribute, or process the signal. This
distinction provides the physical reference for the six analytical families
used in the survey.

The evidence is organized into fiber, free space optical (FSO), visible light
communication and LiFi (VLC/LiFi), photonic terahertz (THz), hybrid optical,
and other optical systems. These analytical categories place each result in
its relevant physical regime before performance comparison.

Fiber systems either sense the guided path directly or use it to transport
signals to a downstream segment that reaches the target. FSO systems place target
geometry and beam control in a directional optical path, while VLC/LiFi adds
illumination and access constraints. In photonic THz systems, optical hardware
may generate, distribute, or process the carrier, but the path to the target
remains radio frequency. Hybrid systems require both domains and their
interface to be retained. These distinctions establish the physical reference
frame. Section~\ref{sec:optical_modality_families} later synthesizes their study
counts, implementations, and constraints specific to each platform
\cite{OISAC_SCR00220,OISAC_SCR00940,OISAC_SCR00196,OISAC_SCR00083,
OISAC_SCR00366}.

\subsection{Integration Mechanisms and Measurement Planes}

Integration is located where the two functions share a component, signal,
resource, path, or decision. Application coupling can coordinate a service even
when the signal paths remain separate. At the physical and signal levels,
systems may share a link, hardware element, optical carrier, waveform, or
allocation rule. Joint design connects these choices by selecting at least one
variable using outcomes from both functions.

These locations overlap within a system. Shared hardware can support separate
waveforms, while a joint optimization model can coordinate both functions
before hardware validation. A shared hardware application
\cite{OISAC_SCR00001}, a fiber system with a shared link
\cite{OISAC_SCR00957}, and a design that combines multiple mechanisms
\cite{OISAC_SCR00220} therefore represent distinct engineering relations. We
retain these attributes separately so later sections can compare concurrency,
resource use, and validation.

The measurement plane adds where the reported quantity is observed. Design
variables describe carrier and waveform settings, launch power, and geometry.
Optical power, loss, and optical signal to noise ratio (OSNR) belong to the
optical link. After conversion, electrical signal to noise ratio (SNR), error
rate, and throughput describe the electrical and digital stages. Sensing
measurements add task specific observations of geometry, motion, environment,
classification, or detection.

Familiar labels become comparable through matched references. OSNR and
electrical SNR belong to different planes, just as launch and received power
describe different locations in the chain. Rate accounting distinguishes gross
from net data, and error accounting distinguishes the event measured before
forward error correction (pre-FEC) from the event measured after it (post-FEC).
Sensing likewise separates range resolution,
estimation error, precision, and analytical bounds. Comparison therefore
follows the task, target, operating point, channel, and processing chain
together with the metric name.

Provenance adds the form of evidence. Analytical bounds, simulations,
laboratory measurements, offline dataset results, and field observations
support different kinds of inference. When decisive conditions align, their
results can reveal a bounded relation across studies. Within a study, a sweep
can show how a design choice changes outcomes across its reported operating
region.

% ============================================================================
% FUTURE FIGURE 2 PRODUCTION SPECIFICATION --- DO NOT MATERIALIZE IN THIS PASS
% Intended placement: here, after integration and measurement context and
% immediately before the formal comparison profile.
%
% Working title:
% Four-Axis O-ISAC Comparison Framework
% Stable label: fig:comparison_framework
% Blueprint ID: FIG-OISAC-01
% Data authority: the Phase D metric, tradeoff, and codebook schemas together
% with the review's comparison-decision rules; conceptual model, no denominator.
%
% Reader question:
% Which contextual axes must be fixed before a reported O-ISAC value can support
% an admissible inference?
%
% Layout: a two-part hub-and-gate diagram. Panel (a) uses four equal axes feeding
% one candidate comparison profile: physical context; coupling location;
% measurement contract; provenance. Physical context includes modality, the
% optical or radio signal path and conversion boundary, both tasks, target,
% geometry, and communication role. Coupling location includes hardware,
% carrier, waveform, resources, link or channel, joint design, and application.
% The measurement contract includes definition, plane, unit, aggregation,
% operating point, channel, target or ground truth, constraints, and baseline.
% Provenance includes value origin, validation setting, exact source locator,
% and configuration continuity. The central profile must carry the warning
% ``missing field = unknown; never infer from a neighboring study.''
%
% Panel (b) sends the candidate profile to three equal outcomes: within-study
% interpretation; conditional cross-study relation; descriptive only. A crossed
% label states that no unconditional cross-platform comparison category exists.
%
% Main visual message:
% Comparison is an evidence decision governed by aligned context and
% provenance, not a consequence of placing two values beside one another.
%
% Future lead-in (activate only after production and audit):
% Figure~\ref{fig:comparison_framework} condenses the review's comparison logic
% into four axes and three admissible uses without turning them into a maturity
% sequence.
%
% Future caption:
% Four-axis framework used to interpret O-ISAC evidence across platforms.
% Physical context, coupling location, measurement contract, and provenance
% feed a candidate profile that supports within-study interpretation, a
% conditional cross-study relation, or descriptive use. Missing fields are not
% inferred, and the framework contains no unconditional cross-platform
% comparison category.
%
% Visual constraints:
% - Figure 2 shows the decision path; Table II summarizes the profile groups.
%   Exact fields and codes remain in S-Data Dictionary and S-Evidence.
% - Give all four axes equal visual weight; do not imply that one modality is a
%   reference standard or that the axes form a sequence.
% - Use explicit outcome labels, not green/red approval symbols.
% - Do not display numerical performance values or an O-ISAC Pareto frontier.
% - Use grayscale-safe shapes and redundant text labels.
% ============================================================================

\begin{table}[!htbp]
\caption{Comparison profile carried with each reported O-ISAC result}
\label{tab:comparison_record}
\centering
\footnotesize
\setlength{\tabcolsep}{3.2pt}
\renewcommand{\arraystretch}{1.04}
\begin{tabularx}{\columnwidth}{@{}>{\raggedright\arraybackslash}p{0.29\columnwidth}
>{\raggedright\arraybackslash}X@{}}
\toprule
\multicolumn{2}{@{}l}{\emph{Comparison profile}} \\
\addlinespace[1pt]
Component & Context carried with the result \\
\midrule
Physical context & Modality, the signal path that reaches the target and its conversion boundary,
communication task or role, sensing task, and target locate the physical regime. \\
Coupling location & A shared component, carrier, waveform, resource, path,
application, or design rule shows where the functions meet. \\
Measurement contract & Metric name and source definition, statistic,
measurement plane, unit and accounting state, operating point, channel,
geometry, environment, reference or ground truth, constraints, and baseline
define the quantity and operating region. \\
Evidence trace & Validation setting, value origin, and exact source locator
connect the result to its evidence. \\
\bottomrule
\end{tabularx}

\vspace{3pt}
\begin{tabularx}{\columnwidth}{@{}>{\raggedright\arraybackslash}p{0.34\columnwidth}
>{\raggedright\arraybackslash}X@{}}
\toprule
\multicolumn{2}{@{}l}{\emph{Analytical use}} \\
\addlinespace[1pt]
Use & Evidence condition \\
\midrule
Interpretation within a study & The result supports a technical interpretation
under the conditions reported by its source. \\
Bounded relation across studies & Decisive fields align and the remaining
differences are stated. \\
Descriptive retention & The record informs mechanism, scope, or feasibility in
its reported context. \\
\bottomrule
\end{tabularx}
\vspace{1pt}
\parbox{\columnwidth}{\footnotesize\emph{Note.} Electronic Supplements
S-Data Dictionary and S-Evidence provide the exact field definitions,
controlled states, coding rules, and provenance for each record. The comparison
decision follows the source fields available for each result.}
\end{table}

\subsection{The Cross Platform Comparison Profile}

The profile becomes comparative when its decisive fields align. A matched
definition and measurement plane establish what is being compared, while the
operating region and evidence trace state where the relation holds. This keeps
each quantitative relation tied to the physical system and the metric's
reported meaning.

One fiber study illustrates this logic. Its communication and localization
outcomes form a bounded relation under shared conditions, while its other
results retain narrower roles within their recorded settings
\cite{OISAC_SCR00056}. The profile makes each role visible and preserves the
platform context for later synthesis.

Throughout the survey, representative studies anchor recurring patterns,
tables summarize the inventory, and prose explains engineering meaning.
Modality and signal path establish physical context, integration locates the
coupling, and the measurement plane identifies the observation. Together with
the evidence trace, these elements determine how evidence can be related
across studies.
